% Generated draft: requires mathematical and source review.
\documentclass[11pt]{article}
\usepackage[margin=25mm]{geometry}
\usepackage{fontspec}
\setmainfont{Latin Modern Roman}
\usepackage{amsmath,amssymb,mathtools,unicode-math}
\setmathfont{Latin Modern Math}
\usepackage{enumitem,needspace,xurl,hyperref}
\setlength{\parindent}{0pt}
\setlength{\parskip}{5pt}
\setlength{\emergencystretch}{3em}
\newcommand{\N}{\mathbb{N}}
\newcommand{\Z}{\mathbb{Z}}
\newcommand{\Q}{\mathbb{Q}}
\newcommand{\R}{\mathbb{R}}
\newcommand{\C}{\mathbb{C}}
\newcommand{\F}{\mathbb{F}}
\newcommand{\A}{\mathbb{A}}
\newcommand{\PP}{\mathbb{P}}
\newcommand{\E}{\mathbb{E}}
\newcommand{\eps}{\varepsilon}
\newcommand{\dd}{\,\mathrm{d}}
\begin{document}
\section*{UnsolvedMath ID: 30003975}
\textbf{Source problem code:} catalogue entry
\subsection{Definitions and mathematical statement}
For the following conventional formulation, let \(R\) be a commutative Noetherian regular ring and let \(I\subsetneq R\) be a radical ideal with finite positive big height
\[
h=\max\{\operatorname{ht}(\mathfrak p):\mathfrak p\in\operatorname{Min}_R(I)\}.
\]
Here \(\operatorname{Min}_R(I)\) denotes the set of prime ideals minimal over \(I\), and regularity means that every localization at a prime ideal is a regular local ring. For an integer \(m\geq1\), define the ordinary and symbolic powers by
\[
I^m=\underbrace{I\cdots I}_{m\text{ factors}},\qquad
I^{(m)}=\bigcap_{\mathfrak p\in\operatorname{Min}_R(I)}\bigl(I^mR_{\mathfrak p}\cap R\bigr),
\]
where the intersection with \(R\) denotes contraction along localization.

The source asks whether a single containment with a fixed offset \(C\) forces the corresponding containments eventually. If \(C\) is interpreted as a fixed integer and only positive symbolic exponents are used, the proposed implication is
\[
\begin{split}
&\left[\exists k\in\Z_{\geq1}\;\bigl(hk-C\geq1\ \text{and}\ I^{(hk-C)}\subseteq I^k\bigr)\right]\\
&\qquad\Longrightarrow
\left[\exists N\in\Z_{\geq1}\;\forall n\in\Z_{\geq N}\;\bigl(hn-C\geq1\ \text{and}\ I^{(hn-C)}\subseteq I^n\bigr)\right].
\end{split}
\]
The same \(R,I,h,C\) occur throughout, and \(N\) may depend on these data. This is a conditional precisification requiring review: the supplied source does not specify the domain, sign, or admissible range of \(C\), any additional lower bound on \(k\), or conventions for nonpositive symbolic exponents. It also does not explicitly specify the ring conventions, properness of \(I\), or exclusion of big height zero. Thus the conventions above do not resolve the intended scope of the original question.
\subsection{Short English statement}
Does one containment \(I^{(hk-C)}\subseteq I^k\), for a radical ideal of big height \(h\), imply \(I^{(hn-C)}\subseteq I^n\) for every sufficiently large integer \(n\), with the same offset \(C\)? The permitted values of \(C\), restrictions on the initial index \(k\), and boundary conventions remain unspecified in the supplied record.
\subsection{Sources}
\begin{itemize}
\item Local UnsolvedMath registry, permanent ID \(30003975\), source code \texttt{catalogue entry}, title ``Stable Symbolic-Power Containments from a Single Containment.'' The supplied complete record is the local registry evidence; project registry locations are \texttt{data/unsolvedmath/problems.json} and the supplementary \texttt{catalog/problems.json}.
\item Supplied source citation: \emph{Mini-Workshop: Asymptotic Invariants of Homogeneous Ideals} (2019), Oberwolfach Reports, DOI 10.4171/owr/2018/45, pp.\ 2739--2767. \url{https://doi.org/10.4171/owr/2018/45}. The record separately gives proposed year 2018 and supplies no named proposer.
\item Supplied literature source, ``Stable Harbourne conjecture'': \url{https://arxiv.org/abs/1809.06955}.
\item Supplied literature source, ``Expected resurgences'': \url{https://arxiv.org/abs/1903.12122}.
\item Literature assessment dated exactly \texttt{2026-08-21}: the supplied curation labels the problem open and states that later papers establish stable containment for important classes and provide sufficient criteria, while the general implication for arbitrary radical ideals remains open. This is supplied dated curation derived from the Oberwolfach report, not an independently verified assessment; its status and cited claims require independent verification.
\item The registry labels the statement \texttt{corrected\_verified}, explaining that mathematical content was preserved while notation, typography, and sentence structure were repaired. This supplied editorial label does not establish mathematical correctness or resolve the missing hypotheses.
\end{itemize}
\end{document}
